\documentclass[11pt,a4paper]{article}

\usepackage[T1]{fontenc}
\usepackage[margin=25mm]{geometry}
\usepackage{microtype}
\usepackage{graphicx}
\usepackage{booktabs}
\usepackage{array}
\usepackage{float}
\usepackage{caption}
\usepackage{enumitem}
\usepackage{xcolor}
\usepackage[hidelinks]{hyperref}
\usepackage{fancyhdr}

\graphicspath{{figures/}}
\definecolor{rulegray}{HTML}{B8C0C8}

\hypersetup{
  pdftitle={AI Programming Foundations Project: Module Summary},
  pdfauthor={Johannes Pawelczyk},
  pdfsubject={Exploratory analysis of the Titanic training dataset as part of the Udacity MSc AI capstone project}
}

\pagestyle{fancy}
\fancyhf{}
\lhead{AI Programming Foundations Project}
\rhead{Module Summary}
\cfoot{\thepage}
\renewcommand{\headrulewidth}{0.4pt}
\setlength{\headheight}{14pt}

\captionsetup{font=small,labelfont=bf}
\setlist{nosep,leftmargin=1.5em}
\setlength{\parindent}{0pt}
\setlength{\parskip}{0.65em}
\renewcommand{\arraystretch}{1.15}

\newcommand{\reportsection}[1]{%
  \section{#1}
  \vspace{-0.4em}
  \textcolor{rulegray}{\rule{\textwidth}{0.5pt}}
  \vspace{0.2em}
}

\newcommand{\reportfigure}[2]{%
  \IfFileExists{figures/#1}{%
    \includegraphics[width=0.92\textwidth]{#1}%
  }{%
    \fbox{\parbox[c][52mm][c]{0.88\textwidth}{\centering\small #2}}%
  }%
}

\title{\textbf{AI Programming Foundations Project}\\[0.4em]
\large Exploratory analysis of the Titanic passenger dataset}
\author{Johannes Pawelczyk}
\date{28. September 2026}

\begin{document}

\maketitle
\thispagestyle{fancy}

\reportsection{Overview}

For this project I used the recommended Titanic training data to build a small,
reproducible analysis in Python. My main research question was how survival differed by
passenger class, sex, age, and family size. I did not train a predictive model; the
purpose was to demonstrate understanding of proper loading, cleaning, summarization,
and plotting on a real dataset.

The submission contains the full analysis. It reads the local CSV file using
Pandas, applies the cleaning functions, produces grouped summaries, and creates
the three figures included below. The repository also contains the source data,
the generated figures, and pinned package versions, as well as brief instructions.

\reportsection{Dataset Description}

I used \emph{train.csv} from Kaggle's
\href{https://www.kaggle.com/c/titanic/data}{Titanic: Machine Learning from
Disaster} competition. It contains 891 passenger records and 12 columns.
The outcome variable, \texttt{Survived}, is binary: 342 passengers are
recorded as survivors and 549 as non-survivors. Thus, the overall recorded survival
rate is 38.4\%.

\begin{table}[H]
  \centering
  \small
  \begin{tabular}{@{}p{0.28\textwidth}p{0.22\textwidth}p{0.40\textwidth}@{}}
    \toprule
    Check & Result & Treatment in the notebook \\
    \midrule
    Missing age & 177 rows (19.9\%) & Median within recorded sex and class; missingness flag retained \\
    Missing cabin & 687 rows (77.1\%) & Raw value removed; cabin-known indicator retained \\
    Missing embarkation port & 2 rows (0.2\%) & Filled with the modal port, S \\
    Exact duplicate rows & 0 & Checked; no rows removed \\
    \bottomrule
  \end{tabular}
  \caption{Data-quality checks performed prior to the analysis.}
  \label{tab:quality}
\end{table}

Important note: The file is a historical passenger list, not a sample designed for
statistical inference. The results below describe patterns in these records. They
should not be interpreted as causal effects.

\reportsection{Workflow Description}

The notebook is broken up into six stages:

\begin{enumerate}
  \item Import required libraries and set display and plotting options.
  \item Load CSV and inspect its rows, dimensions, data types, and missingness.
  \item Clean column names/values, handle missing data, and create derived fields.
  \item Calculate counts and survival rates for the groups used in the analysis.
  \item Plot the main comparisons with titles, axis labels, and group sizes.
  \item Record findings, assumptions, and limitations.
\end{enumerate}

Two functions perform the cleaning step. \texttt{normalize\_column\_names} converts
the original headings to lower-case snake case. \texttt{clean\_titanic\_data}
then handles duplicates, whitespace, missing ages, embarkation values, cabin
availability, and the derived family-size fields. A third function,
\texttt{summarize\_survival}, returns the passenger count, survivor count, and
survival rate for any requested grouping.

\reportsection{Key Decisions \& Assumptions}

All 891 rows were kept, as dropping every incomplete row would discard most of the
dataset because cabin is missing for 77.1\% of passengers. Instead, I replaced
the raw cabin field with a \texttt{cabin\_known} indicator. This preserves the
fact that a cabin was recorded without treating an incomplete cabin string as a
reliable location.

Age has less missingness but is still absent for 177 passengers. I imputed those
values as the median for the passenger's recorded sex and class. Furthermore, the
separate \texttt{age\_missing} flag distinguishes observed and imputed
values. Finally, the two missing embarkation values were filled with the most common
port.

Family size is defined as \texttt{SibSp + Parch + 1}. This counts the passenger,
their siblings/spouse, and their listed parents or children. It does not capture
companions outside those relationships, so it is only a proxy for the actual travel
group. The cleaned table otherwise follows a consistent structure, with each variable
having its own column and each passenger having one row \cite{wickham2014}.

\reportsection{Results \& Interpretation}

The largest differences appear when survival is grouped by passenger class and sex.
Female passengers have a higher survival rate in every class. The rate ranges from
96.8\% for first-class women to 13.5\% for third-class men. Class still matters within
each recorded sex category.

\begin{figure}[H]
  \centering
  \reportfigure{figure_1_survival_by_class_and_sex.png}{Survival rate by passenger class and sex}
  \caption{Recorded survival rate by passenger class and sex.}
  \label{fig:class-sex}
\end{figure}

The age distributions overlap substantially. Both outcome groups have a raw
median age of approximately 28 years, although the survivor distribution has a visible
peak among children. Because 177 ages were filled before plotting, small local
differences should be interpreted cautiously.

\begin{figure}[H]
  \centering
  \reportfigure{figure_2_age_by_survival.png}{Age distribution by survival status}
  \caption{Age distributions for recorded survivors and non-survivors.}
  \label{fig:age}
\end{figure}

Passengers in family groups of two to four have higher recorded survival rates
than passengers travelling alone. The relationship is not linear: the rate
drops for groups of five or more. Several large-family categories contain very
few passengers. For example, family sizes 8 and 11 contain only 6 and 7 records,
so their rates are unstable and should not support broad conclusions.

\begin{figure}[H]
  \centering
  \reportfigure{figure_3_survival_by_family_size.png}{Figure preview: survival rate by family size}
  \caption{Recorded survival rate by derived family size; labels show group counts.}
  \label{fig:family}
\end{figure}

\reportsection{Responsible Practice: Bias \& Data Quality}

Cleaning choices can change the analysis results and conclusions. A
complete-case approach would remove 19.9\% of rows because of age alone, while
requiring cabin would remove 77.1\%. If missingness differs by class or recorded
sex, those deletions would also change the comparisons. I therefore kept the
rows, recorded missingness explicitly, and reported group counts next to rates.

However, median imputation has its own limitations. It reduces variation
and can reinforce the group structure used to calculate the median. A follow-up
analysis should compare the current results with complete-case estimates and other
imputation methods. The original categories also reflect the conventions and unequal
social conditions of the time period. Any predictive use would need separate subgroup
checks, a clear reason for using sensitive attributes, and safeguards against
data leakage.

\reportsection{Reproducibility}

The notebook uses relative paths and runs from the repository root. All code
cells have been executed without errors. Assertions check the row count, required
columns, survival labels, filled age and embarkation values, and the range of the
derived family-size field. Package versions are listed in \texttt{requirements.txt}
and the README explains how to create the environment and rerun the notebook.

The Git history separates the initial project structure, the executed analysis,
and the report. The analysis was developed on a \texttt{development} branch and
merged into \texttt{main}. Keeping the notebook, narrative, environment details,
and source data together follows the reproducible teaching approach described by
Danchev \cite{danchev2022}.

\reportsection{Sources \& Citations}

\begin{thebibliography}{9}

\bibitem{danchev2022}
Danchev, V. (2022). Reproducible data science with Python: An open learning
resource. \emph{Journal of Open Source Education, 5}(56), 156.
\url{https://doi.org/10.21105/jose.00156}

\bibitem{kaggle}
Kaggle. (n.d.). \emph{Titanic: Machine Learning from Disaster -- Data}.
Retrieved 28 September 2026 from
\url{https://www.kaggle.com/c/titanic/data}

\bibitem{wickham2014}
Wickham, H. (2014). Tidy data. \emph{Journal of Statistical Software, 59}(10),
1--23. \url{https://doi.org/10.18637/jss.v059.i10}

\end{thebibliography}

\end{document}
